\corrigesChapitre{11}{Poids, masse et gravitation}
\par\noindent\begin{minipage}{\linewidth}
\begin{corrige}{1}{Associer les unités}
\begin{questions}
\item Kilogramme pour la masse ; newton pour le poids.
\item Le dynamomètre.
\end{questions}
\end{corrige}
\end{minipage}
\par\noindent\begin{minipage}{\linewidth}
\begin{corrige}{2}{Le lieu change}
\begin{questions}
\item Non, s’il ne gagne ni ne perd de matière.
\item Oui, car le coefficient de pesanteur local change.
\end{questions}
\end{corrige}
\end{minipage}
\par\noindent\begin{minipage}{\linewidth}
\begin{corrige}{3}{Convertir une masse}
\begin{questions}
\item $0{,}250\unite{kg}$.
\item $1\,500\unite g$.
\end{questions}
\end{corrige}
\end{minipage}
\par\noindent\begin{minipage}{\linewidth}
\begin{corrige}{4}{Un poids simple}
\begin{questions}
\item $P=2\times10=20\unite N$.
\item Direction verticale ; sens vers le bas, vers le centre de la Terre dans le modèle.
\end{questions}
\pcFigureEx{figures/cor-poids-20N.pdf}{Poids de 20 N, vertical vers le bas ; schéma qualitatif sans échelle de longueur.}
\end{corrige}
\end{minipage}
\par\noindent\begin{minipage}{\linewidth}
\begin{corrige}{5}{Une interaction}
\begin{questions}
\item Non : elle s’exerce à distance.
\item Oui : les objets massifs s’attirent.
\end{questions}
\end{corrige}
\end{minipage}
\par\noindent\begin{minipage}{\linewidth}
\begin{corrige}{6}{Un petit objet}
\begin{questions}
\item $80\unite g=0{,}080\unite{kg}$ ; $P=0{,}80\unite N$.
\item On trouverait à tort $800\unite N$, soit mille fois trop.
\end{questions}
\par\smallskip\noindent\textbf{Pour comprendre.}\enspace Convertir d’abord $80\unite g=0{,}080\unite{kg}$. Alors $P=mg=0{,}080\times10=0{,}80\unite N$. Employer 80 à la place de 0,080 produirait une valeur 1 000 fois trop grande.\par
\end{corrige}
\end{minipage}
\par\noindent\begin{minipage}{\linewidth}
\begin{corrige}{7}{Retrouver la masse}
\begin{questions}
\item $m=P/g=35/10=3{,}5\unite{kg}$.
\item $3\,500\unite g$.
\end{questions}
\par\smallskip\noindent\textbf{Pour comprendre.}\enspace De $P=mg$, déduire $m=P/g=35/10=3{,}5\unite{kg}$ ; puis convertir en 3 500 g. Le poids en N n’est pas une masse en kg.\par
\end{corrige}
\end{minipage}
\par\noindent\begin{minipage}{\linewidth}
\begin{corrige}{8}{Estimer le coefficient}
\begin{questions}
\item $g=P/m=4{,}9/0{,}50=9{,}8\unite{N/kg}$.
\item Newton par kilogramme.
\end{questions}
\par\smallskip\noindent\textbf{Pour comprendre.}\enspace La relation donne $g=P/m=4{,}9/0{,}50=9{,}8\unite{N/kg}$. Le quotient compare des newtons à des kilogrammes ; son unité n’est pas N seule.\par
\end{corrige}
\end{minipage}
\par\noindent\begin{minipage}{\linewidth}
\begin{corrige}{9}{Terre et Lune}
\begin{questions}
\item $P_T=50\unite N$ ; $P_L=8\unite N$.
\item $50/8=6{,}25$ ; le poids terrestre est environ six fois plus grand avec ces valeurs arrondies.
\end{questions}
\pcFigureEx{figures/cor-terre-lune-5kg.pdf}{Même masse de 5 kg ; poids de 50 N et 8 N. Les longueurs des flèches respectent le même rapport.}
\end{corrige}
\end{minipage}
\par\noindent\begin{minipage}{\linewidth}
\begin{corrige}{10}{Une droite}
\begin{questions}
\item Une proportionnalité de coefficient $10\unite{N/kg}$.
\item $P=0{,}45\times10=4{,}5\unite N$.
\end{questions}
\par\smallskip\noindent\textbf{Pour comprendre.}\enspace Les trois quotients P/m valent 10 N/kg. Le modèle est donc $P=10m$ dans la plage étudiée ; pour 0,45 kg, on obtient 4,5 N.\par
\pcFigureEx{figures/cor-poids-graphe.pdf}{Le quotient P/m vaut 10 N/kg ; le point prévu pour 0,45 kg est sur la même droite.}
\end{corrige}
\end{minipage}
\par\noindent\begin{minipage}{\linewidth}
\begin{corrige}{11}{Une balance détournée}
\begin{questions}
\item Son réglage suppose la pesanteur terrestre ; une même masse exerce une force plus faible sur la Lune.
\item Une balance comparant l’objet à des masses étalons dans les mêmes conditions, si son fonctionnement est adapté.
\end{questions}
\end{corrige}
\end{minipage}
\par\noindent\begin{minipage}{\linewidth}
\begin{corrige}{12}{La phrase du quotidien}
\begin{questions}
\item Sa masse est de $50\unite{kg}$.
\item Son poids vaut $P=50\times10=500\unite N$ dans ces conditions.
\end{questions}
\end{corrige}
\end{minipage}
\par\noindent\begin{minipage}{\linewidth}
\begin{corrige}{13}{Un zéro décalé}
\begin{questions}
\item Un dynamomètre mal réglé au zéro peut produire ce décalage.
\item Non : on contrôle d’abord le zéro, le protocole, les unités et les autres sources d’erreur.
\end{questions}
\pcFigureEx{figures/cor-poids-decale.pdf}{Un décalage constant de 0,3 N peut provenir de l’appareil ; contrôler le zéro avant de modifier le modèle.}
\end{corrige}
\end{minipage}
\par\noindent\begin{minipage}{\linewidth}
\begin{corrige}{14}{Une orbite}
\begin{questions}
\item Non : l’attraction terrestre contribue au mouvement orbital.
\item La station et les objets qu’elle contient sont en chute commune autour de la Terre ; l’impesanteur apparente ne signifie pas absence d’interaction.
\end{questions}
\end{corrige}
\end{minipage}
\par\noindent\begin{minipage}{\linewidth}
\begin{corrige}{15}{Une donnée de gravitation}
\begin{questions}
\item $12/4=3\unite N$.
\item Non : il faut connaître le système, les masses, les distances et le domaine du modèle. Ici on exploite seulement une règle fournie pour préparer la suite.
\end{questions}
\end{corrige}
\end{minipage}
\par\noindent\begin{minipage}{\linewidth}
\phantomsection\label{cor-defi-11}\pcSolutionDefi{11}{Le colis lunaire}\begin{questions}
\item $m=0{,}750\unite{kg}$ ; $P_T=7{,}35\unite N$ et $P_L=1{,}20\unite N$.
\item Le calibre 10 N convient ; 1 N et 5 N seraient dépassés.
\item On ne conclut pas à une masse nulle. L’instrument en chute ne fonctionne plus comme une suspension statique ; la masse reste $0{,}750\unite{kg}$.
\end{questions}
\pcFigureEx{figures/cor-colis-lunaire.pdf}{Comparer deux poids ne change pas la masse ; les calibres doivent couvrir la force mesurée.}
\end{minipage}
\par\noindent\begin{minipage}{\linewidth}
\phantomsection\label{cor-activite-11}\pcSousTitre{Pistes pour l’investigation}
Convertir les masses en kg puis tracer P en N en fonction de m. Une droite près de l’origine et des quotients proches de $9{,}8$ à $10\unite{N/kg}$ sont attendus sur Terre dans ces conditions. La valeur retenue doit respecter la précision des masses et du dynamomètre. Une ordonnée à l’origine non nulle peut signaler un décalage du zéro.\par\end{minipage}
